\documentclass[../main.tex]{subfiles}

\begin{document}
\addcontentsline{toc}{chapter}{Preface}
\chapter*{Preface}

\addcontentsline{toc}{section}{About the Book}
\section*{\S\! About the Book}

Teaching and sharing information, in my opinion, is the best way to
learn. Whenever I explain the contents that I learned to others,
I start to really distinguish what I know and what I do not know.
Therefore when writing notes, I write it as if I am explaining it to
someone, sharing my thought process and understandings.

LibreNotebook is a collection of my personal notes. To be honest,
the contents are biased towards my interests, and includes topics
that I am passionate about and concepts that I really wish to
learn thoroughly. Having said that, writing requires time and
dedication. I write most of my notes on physical notebook, and this
includes the selected topics that I dedicated to type. Moreover,
it is not meant to replace textbooks, courses, and dedicated
resources. On the other hand, if you are also interested and learning
more about the topics included in my notes, I believe LibreNotebook
can serve as a great resource to see how other person understood
the concept.

\addcontentsline{toc}{section}{Purpose and Motivation}
\section*{\S\! Purpose and Motivation}

LibreNotebook initially started with my personal notes. To be exact,
it began with me studying for national math olympiad. Frankly,
I was lost when studying for the national math olympiad as not
many resources were available specifically for the high school
division of our country's olympiad. Even if there are resources,
it was either for middle school round or was commercial with high
expense. Naturally, I began seeking for problems and resources from
different countries. There, with open and public resources, I was able
to learn and improve. Since then, I always wanted to return the
blessings received from open education.

Over the 2026 summer break, I wrote notes on linear algebra and
machine learning. With increasing lengths and details, I thought they
could be more than just personal notes. With different handouts that I
wrote for multiple purposes, I thought it would be possible to return
the blessings received by doing the same of returning knowledge.
LibreNotebook is the notebook that highlights understandings that were
learned to be shared.

\addcontentsline{toc}{section}{Contributing}
\section*{\S\! Contributing}

First and foremost, thank you for your interest and taking your time
to contribute. Any contributions ranging from a typo to writing a new
note is appreciated! The source files for LibreNotebook are hosted
in \href{https://github.com/enochyu-official/LibreNotebook}{GitHub},
and you can contribute by creating an issue, submitting a pull
request, or write a mail using the contact form from
\href{https://enochyu.com#contact}{my website}.

Although contributions are welcomed, I want to include a notice on
AI usage. I know this topic is philosophical and political,
but I want to say I really don't want LibreNotebook to be flooded with
AI generated contents. Unlike AI assisted contents, many AI generated
contents lack proper credits and quality checks. For transparency of
contributors, it would be great if AI usage is disclosed for each
contribution. With that, thank you for your interest and efforts!

\addcontentsline{toc}{section}{About the Author}
\section*{\S\! About the Author}

Enoch is an ardent student who retains keen interest in pure math,
machine learning, and quantitative finance. His previous
accomplishments includes honorable mention from Korean Mathematical
Olympiad (~300 in South Korea) and numerous awards from AMC.
He currently dedicates his time on making
\href{https://www.youtube.com/@enochyu_official}{YouTube videos}
on various math topics and
\href{https://nanumlearn.enochyu.com}{Nanum Learn}. More on Enoch can
be found in his \href{https://enochyu.com}{personal website}.

\end{document}

